%%
%% Dynamic Platoon Formation of Multi-Type Autonomous Vehicles
%% for Sustainable Urban Mobility
%%
%% ACM SIGSPATIAL 2026 Short Paper Track Submission
%% Single-blind, 4 pages including references
%%
%% Commands for TeXCount
%TC:macro \cite [option:text,text]
%TC:macro \citep [option:text,text]
%TC:macro \citet [option:text,text]
%TC:envir table 0 1
%TC:envir table* 0 1
%TC:envir tabular [ignore] word
%TC:envir displaymath 0 word
%TC:envir math 0 word
%TC:envir comment 0 0

\documentclass[sigconf]{acmart}

%% Additional packages used in our paper
\usepackage{algorithm}
\usepackage{algpseudocode}
\usepackage{tikz}
\usetikzlibrary{patterns,arrows.meta,positioning,calc}

\AtBeginDocument{%
  \providecommand\BibTeX{{%
    Bib\TeX}}}

%% ---------------------------------------------------------------------------
%% Rights management
%% These will be replaced by the official values from ACM rights confirmation
%% after acceptance. SAMPLE values for now.
%% ---------------------------------------------------------------------------
\setcopyright{acmlicensed}
\copyrightyear{2026}
\acmYear{2026}
\acmDOI{XXXXXXX.XXXXXXX}
\acmConference[SIGSPATIAL '26]{The 34th ACM SIGSPATIAL International Conference on Advances in Geographic Information Systems}{November 3--6, 2026}{Riverside, CA, USA}
\acmISBN{978-1-4503-XXXX-X/26/11}

%% EasyChair submission ID
\acmSubmissionID{83}

%%
%% End of preamble.
\begin{document}

%% ===========================================================================
%% TITLE
%% ===========================================================================
\title{Dynamic Platoon Formation of Multi-Type Autonomous Vehicles for Sustainable Urban Mobility}

%% ===========================================================================
%% AUTHORS (single-blind: names required)
%% ===========================================================================
\author{Jaeyun Ree}
\email{jree0010@student.monash.edu}
\affiliation{%
  \institution{Monash University}
  \city{Melbourne}
  \state{VIC}
  \country{Australia}
}

\author{Mohammed Eunus Ali}
\email{eunus.ali@monash.edu}
\affiliation{%
  \institution{Monash University}
  \city{Melbourne}
  \state{VIC}
  \country{Australia}
}

\author{Muhammad Aamir Cheema}
\email{aamir.cheema@monash.edu}
\affiliation{%
  \institution{Monash University}
  \city{Melbourne}
  \state{VIC}
  \country{Australia}
}

\renewcommand{\shortauthors}{Ree, Ali, and Cheema}

%% ===========================================================================
%% ABSTRACT (from ITSC, lightly trimmed)
%% ===========================================================================
\begin{abstract}
This paper addresses the energy inefficiency of single-occupancy vehicles by introducing a novel cooperative autonomous vehicle system where smaller Passive Vehicles (PVs) can be physically towed by larger Active Vehicles (AVs) in a shared road segment. Unlike traditional platooning that relies solely on aerodynamic drafting, our approach enables PV propulsion energy to be fully offset by the towing AV during the shared segment travel, though the AV bears the additional load. We formulate the platoon formation problem as a constrained optimization problem and propose two algorithms: a Greedy Maximum-Weight Matching algorithm and an Iterative Linear Assignment (ILA) algorithm using the Jonker--Volgenant method. A key innovation is the multi-segment matching capability, allowing a single PV to be towed by multiple AVs across different route segments. Experimental evaluation on synthetic scenarios with up to 400 PVs and 80 AVs demonstrates that both algorithms achieve 25--52\% distance coverage depending on capacity, with ILA consistently matching or exceeding the greedy baseline while running 2.8--11.5$\times$ faster.
\end{abstract}

%% ===========================================================================
%% CCS CONCEPTS
%% TODO: Generate the actual XML at https://dl.acm.org/ccs/ccs.cfm
%% Suggested categories to select on that site:
%%   - Theory of computation > Design and analysis of algorithms >
%%     Mathematical optimization > Combinatorial optimization
%%   - Applied computing > Operations research > Transportation
%%   - Computing methodologies > Modeling and simulation
%% Paste the generated <ccs2012>...</ccs2012> block here, then list ccsdesc lines.
%% ===========================================================================
\begin{CCSXML}
<ccs2012>
   <concept>
       <concept_id>TODO</concept_id>
       <concept_desc>TODO: replace with CCS-generated concept</concept_desc>
       <concept_significance>500</concept_significance>
   </concept>
</ccs2012>
\end{CCSXML}

\ccsdesc[500]{TODO: replace with CCS-generated concept~e.g. Theory of computation~Combinatorial optimization}

%% ===========================================================================
%% KEYWORDS (must match EasyChair submission)
%% ===========================================================================
\keywords{autonomous vehicles, vehicle platooning, spatial-temporal matching, linear assignment, sustainable mobility}

\maketitle

%% ===========================================================================
%% SECTION 1: INTRODUCTION
%% Target length: ~1.2 of 4 pages.
%% This is the most important section to rewrite given ITSC reviewer feedback.
%% Eunus's meeting priorities to address here:
%%   1. Explicit motivation that towing-based AV-PV system is a plausible
%%      near-future scenario (cite vision papers / patents / autonomous-driving
%%      roadmaps).
%%   2. Explicit SCOPE statement: we focus on the algorithmic / spatial-
%%      temporal matching problem; mechanical, safety, and regulatory aspects
%%      are orthogonal and out of scope.
%%   3. Re-frame the contribution toward SIGSPATIAL's spatial-algorithm
%%      audience (away from ITSC's mechanical-feasibility framing).
%%   4. Fold the deleted Related Work section into ONE paragraph here.
%%
%% Content to migrate from conference_101719.tex:
%%   - L. 74        : single-occupancy / GHG motivation
%%   - L. 131       : near-future AV coordination paragraph
%%   - L. 133       : problem statement paragraph
%%   - L. 136       : Figure 1 walk-through (keep the figure)
%%   - L. 138-140   : prior work (compress into 1 paragraph)
%%   - L. 142       : problem statement + algorithm overview
%%   - L. 144       : contributions list
%% ===========================================================================
\section{Introduction}
\label{sec:introduction}

%% TODO: Paragraph 1 -- Motivation (energy / mobility crisis).
%% Source: ITSC L.74. Keep ~3 sentences. Cite usdot2024, epa2023.

%% TODO: Paragraph 2 -- Near-future AV coordination is plausible.
%% This is the NEW paragraph addressing ITSC reviewer feedback.
%% Cite (a) vision papers on cooperative/connected autonomous vehicles,
%% (b) any white paper or patent that proposes physical coupling between road
%%     vehicles (mention train-style modular concepts as a known analogue),
%% (c) yurtsever2020survey, shladover2009cooperative as warm-ups.
%% Eunus on 2026-05-08: explicitly say that mechanical, safety, regulatory
%% details are orthogonal and OUT OF SCOPE; we focus on the algorithmic side.

%% TODO: Paragraph 3 -- Introduce AV / PV terminology and the dynamic
%% platoon formation problem. Source: ITSC L.133. Keep tight.

%% TODO: Paragraph 4 -- Figure 1 walk-through (illustrative example).
%% Source: ITSC L.136. Keep figure; trim text to 2--3 sentences.
%% Figure migration: copy the tikzpicture from ITSC L.78--123 verbatim;
%% ensure it fits one column in acmart layout.

\begin{figure}[t]
\centering
%% TODO: paste TikZ figure from conference_101719.tex (L.78--123).
%% After paste, verify caption matches ACM style and fits one column.
\caption{Conceptual overview of towing-based cooperation. Activators $A_1$ and $A_2$ each have capacity 2; $T2$ is towed by $A_1$ to $x{=}60$, self-drives to $x{=}70$, then is handed off to $A_2$.}
\Description{A one-dimensional highway with three passive vehicles and two active vehicles, showing dashed self-driving segments and solid towed segments along positions 0 to 100.}
\label{fig:intro_concept}
\end{figure}

%% TODO: Paragraph 5 -- Prior work compressed into ONE paragraph.
%% Source: ITSC L.138, L.140, L.142 + the old Related Work section
%% (L.149--158 in ITSC). Three threads to mention briefly:
%%   (i)   Platooning / drafting (cite tsugawa2016review, bergenhem2012overview).
%%   (ii)  Ridesharing / assignment (cite furuhata2013ridesharing,
%%         agatz2012optimization, alonso2017demand).
%%   (iii) Combinatorial assignment algorithms (cite kuhn1955hungarian,
%%         munkres1957algorithms, cattrysse1992survey).
%% Make the gap explicit: none of these capture (a) sustained physical
%% coupling, (b) multi-segment coverage of one PV by multiple AVs, and
%% (c) point-wise AV capacity constraints simultaneously.

%% TODO: Paragraph 6 -- Approach overview (Greedy + ILA), framed for
%% SIGSPATIAL's spatial-algorithm audience. Source: ITSC L.142.
%% Emphasize: spatial-temporal matching, bipartite subproblems via
%% Jonker--Volgenant (scipy2020), virtual-slot capacity encoding.
%% Define LSAP and GAP here (reviewer 3 complaint).

%% TODO: Paragraph 7 -- Contributions (3 bullets or 3 sentences).
%% Source: ITSC L.144.
%%   1. Formalize dynamic platoon formation with multi-segment matching
%%      and point-wise capacity, on a 1D highway abstraction.
%%   2. Two solution strategies: greedy + ILA.
%%   3. Controlled experimental study showing 25--52% coverage with ILA
%%      matching or exceeding greedy at 2.8--11.5x lower runtime.
%% Frame 2D extension as future work to pre-empt the obvious reviewer push.


%% ===========================================================================
%% SECTION 2: PROBLEM FORMULATION
%% Target length: ~1 page.
%% Strategy: keep notation table compact, merge System Model + Shared Path
%% + Multi-Segment + Optimization. Drop Illustrative Example subsection
%% (already covered by Figure 1 in Introduction).
%%
%% Content to migrate from conference_101719.tex:
%%   - L. 166-194  : Notation Summary table (KEEP, with units added)
%%   - L. 196-206  : System Model + Active/Passive Vehicle definitions
%%   - L. 208-217  : Shared Path and Energy Model
%%   - L. 219-228  : Multi-Segment Matching definition
%%   - L. 230-250  : Optimization Problem formulation
%%   - L. 252-263  : Assumptions and Temporal Constraints (KEEP, this
%%                   addresses several reviewer concerns)
%% ===========================================================================
\section{Problem Formulation}
\label{sec:formulation}

%% TODO: Lead paragraph -- 1D highway abstraction is a deliberate
%% modeling choice to isolate the matching problem; 2D extension is
%% future work. Reviewer 1 + 3 of ITSC pushed back here, so be explicit.

\subsection{Notation and System Model}
%% TODO: Compact notation table. Source: ITSC L.170--194.
%% IMPORTANT FIX (Reviewer 3, ITSC): add UNITS for every dimensional
%% quantity in the description column:
%%   L       -> meters
%%   e^a, x^a, e^p, x^p, cp, dp, ell  -> meters (position)
%%   C_i     -> integer (vehicles)
%%   t_*, hat_t   -> seconds
%%   v_*     -> meters per second
%%   d_ij, S_ij   -> meters (saving proxy)
%%   L_min   -> meters
%%   tau     -> seconds
%% Also explicitly say capacity is a scalar per AV, not a vector,
%% to fix Reviewer 3's confusion between body text and Table II.

%% TODO: One paragraph definition of Active Vehicle and Passive Vehicle
%% (source: ITSC L.200--206). Trim to essentials.

\subsection{Shared Path and Energy Saving}
%% TODO: Define shared travel distance d_ij and saving S_ij = d_ij when
%% d_ij >= L_min, else 0. Source: ITSC L.208--217.
%% Add one sentence: "We use towed distance as a proxy for PV-side
%% propulsion energy reduction and reserve physics-level modeling for
%% future work" (already in ITSC abstract; restate here for review3).

\subsection{Multi-Segment Matching and Point-wise Capacity}
%% TODO: Define multi-segment matching M_j (a PV can be towed by
%% multiple AVs across disjoint route segments). Source: ITSC L.219--228.
%% State point-wise capacity constraint formally: for every AV a_i and
%% every position x along its route, |{p : (a_i, p) active at x}| <= C_i.
%% This is the spatial-temporal constraint that distinguishes our
%% formulation from standard GAP / LSAP -- emphasise for SIGSPATIAL.

\subsection{Optimization Problem}
%% TODO: Compact problem statement -- maximise total towed distance
%% subject to (a) overlap >= L_min, (b) point-wise capacity, (c)
%% temporal synchronisation. Source: ITSC L.230--250.
%% Spell out LSAP (Linear Sum Assignment Problem) and GAP (Generalized
%% Assignment Problem) on first use, with citations
%% (kuhn1955hungarian, cattrysse1992survey).

\subsection{Assumptions}
%% TODO: BULLETIZED assumption list. Source: ITSC L.252--263.
%% Critical to address ITSC reviewer pushback head-on:
%%   - Known, fixed routes; AVs do not re-route.
%%   - Constant per-vehicle speed; deterministic entry times.
%%   - Coupling/decoupling cost negligible at the algorithmic layer
%%     (reviewer 1 + 3 of ITSC complained -- acknowledge as a modeling
%%     simplification and cite as future work).
%%   - No induced congestion: total vehicle count is unchanged; only
%%     who tows whom changes (reviewer 1 of ITSC).
%%   - Mechanical, safety, and regulatory aspects are out of scope
%%     (refer back to Section 1 scope statement).


%% ===========================================================================
%% SECTION 3: PROPOSED ALGORITHMS
%% Target length: ~0.8 page.
%% Strategy: ONE pseudocode block (ILA only). Greedy described in prose
%% in 2--3 sentences. Drop Theoretical Analysis subsection; fold one
%% sentence about per-iteration optimality into the ILA description.
%%
%% Content to migrate from conference_101719.tex:
%%   - L. 382-419  : Greedy -- compress to 1 short paragraph + complexity
%%   - L. 421-517  : ILA -- KEEP the pseudocode (renamed to algpseudocode),
%%                   trim explanatory text
%%   - L. 519-544  : Theoretical Analysis -- collapse to 1 sentence
%%                   ("per-iteration optimality of the bipartite subproblem
%%                    is guaranteed by the Jonker-Volgenant solver")
%% ===========================================================================
\section{Proposed Algorithms}
\label{sec:algorithms}

\subsection{Greedy Maximum-Weight Matching}
%% TODO: 1 short paragraph -- generate all feasible (AV,PV,segment)
%% candidates, sort by saving, commit greedily while respecting point-
%% wise capacity. Mention worst-case O((NM)^2 log NM). Source: ITSC
%% L.382--419. Keep tight; this is the BASELINE.

\subsection{Iterative Linear Assignment (ILA)}
%% TODO: 2--3 paragraph explanation of ILA, then ONE pseudocode block.
%% Source: ITSC L.421--517.
%% Key points to keep:
%%   - Each iteration: build bipartite graph (PVs vs AV virtual slots),
%%     solve LSAP via Jonker--Volgenant (scipy2020).
%%   - Virtual-slot encoding handles point-wise capacity inside the LSAP.
%%   - Update PV state after each iteration; repeat until no feasible
%%     match remains.
%%   - Per-iteration optimality on the linear subproblem (one sentence
%%     replacing the dropped Theoretical Analysis subsection).
%% Migrate pseudocode using algpseudocode (NOT algorithmic) to avoid
%% acmart compatibility warnings.

%% TODO: Migrate pseudocode here. Switch \algorithmic -> algpseudocode.
\begin{algorithm}
\caption{Iterative Linear Assignment (ILA)}
\label{alg:ila}
\begin{algorithmic}[1]
%% TODO: paste algorithm body from ITSC, rewritten with algpseudocode.
\end{algorithmic}
\end{algorithm}


%% ===========================================================================
%% SECTION 4: EXPERIMENTAL EVALUATION
%% Target length: ~0.6 page.
%% Strategy: keep ONLY the two strongest sweeps -- capacity and length.
%% Drop PV/AV count sweep (or mention in 1 sentence). Drop Discussion
%% and Limitations subsections; fold key points into Results.
%%
%% Content to migrate from conference_101719.tex:
%%   - L. 549-591  : Experimental Setup -- compress to 1 short paragraph
%%   - L. 593-693  : Results -- keep figures for capacity + length
%%                   sweeps; drop PV/AV sweep figure if room is tight.
%%   - L. 695-705  : Discussion -- fold 2--3 sentences into Results.
%%   - L. 707-712  : Limitations -- fold one sentence into Conclusion.
%%
%% Figures to copy from paper/figures/ to paper/sigspatial/figures/:
%%   - capacity_baseline_comparison.png
%%   - capacity_saving_percent_comparison.png  (or _relative_performance)
%%   - length_saving_percent_comparison.png    (or _relative_performance)
%%
%% IMPORTANT FIX (Reviewer 3, ITSC):
%%   - Figure 7/8: clarify difference between "total saving" and
%%     "saving %", or drop one to avoid confusion.
%%   - State unit of L explicitly (meters), and clarify whether L=1600
%%     means 1.6 km (not 1600 km).
%%   - Define density (low/mid/high) and how it varies in the sweeps.
%% ===========================================================================
\section{Experimental Evaluation}
\label{sec:experiments}

\subsection{Setup}
%% TODO: Synthetic data generator, fleet sizes (50--80 AVs, 200--400
%% PVs), capacity range, highway length range with explicit units.
%% Source: ITSC L.549--591. Compress aggressively.

\subsection{Results}
%% TODO: Two paragraphs + two figures (capacity sweep, length sweep).
%% Headline numbers: 25--52% saving range, ILA 0.2--0.9 pp higher than
%% greedy, 2.8--11.5x faster.
%% Fold 1--2 sentences of Discussion (when does ILA matter most) and
%% 1 sentence on Limitations (1D abstraction) into the closing paragraph.


%% ===========================================================================
%% SECTION 5: CONCLUSION
%% Target length: ~0.2 page (3--4 sentences).
%% Strategy: combine Conclusion + Future Work in one short section.
%% State the 2D extension as the headline future direction (this is
%% what we promised Eunus to defer to the summer project).
%% ===========================================================================
\section{Conclusion and Future Work}
\label{sec:conclusion}

%% TODO: 3--4 sentences:
%%   - Summarise the dynamic platoon formation problem + ILA result.
%%   - 1D abstraction is a deliberate modeling choice; extension to 2D
%%     road networks with routing decisions is the primary future work
%%     (pre-empts the obvious reviewer ask).
%%   - Briefly mention: relaxing constant-speed and deterministic-entry
%%     assumptions, modelling coupling/decoupling overhead, energy model
%%     refinement.


%% ===========================================================================
%% ACKNOWLEDGMENTS
%% Use the acks environment (NOT \section{Acknowledgments}).
%% ===========================================================================
%% TODO: decide whether to include acks for short paper. If yes:
%% \begin{acks}
%%   ...
%% \end{acks}


%% ===========================================================================
%% BIBLIOGRAPHY
%% Inline thebibliography for now (same approach as ITSC version).
%% Switch to BibTeX (\bibliography{references}) later if time permits.
%% ===========================================================================
\begin{thebibliography}{00}
%% TODO: migrate \bibitem entries from conference_101719.tex L.727+.
%% Add NEW references for:
%%   - Vision papers / patents on cooperative AV coupling (motivation
%%     section, paragraph 2).
%%   - Train-style modular vehicle concept reference (analogue).
%%   - SIGSPATIAL-aligned spatial-temporal matching references if useful.
\end{thebibliography}

\end{document}
